\begin{solution}
\[
\nabla_\mu J^\mu
=(\nabla_\mu T^{\mu\nu})\xi_\nu
+T^{\mu\nu}\nabla_\mu\xi_\nu=0,
\]
because the first term vanishes and the second contracts a symmetric tensor
with \(\nabla_{(\mu}\xi_{\nu)}=0\).  Stokes' theorem makes
\(\int_\Sigma J^\mu\dd\Sigma_\mu\) hypersurface independent when no flux
crosses the timelike boundary and the relevant integrals converge.
\end{solution}

\begin{solution}
For
\[
T_{\rm EM}^{\mu\nu}
=F^{\mu\rho}F^\nu{}_\rho-\frac14g^{\mu\nu}F_{\rho\sigma}F^{\rho\sigma},
\]
differentiate, use \(\nabla_\mu F^{\mu\rho}=J^\rho\), and contract the Bianchi
identity \(\nabla_{[\mu}F_{\rho\sigma]}=0\).  The derivative terms combine to
\(-F^\nu{}_\rho J^\rho\), with the sign fixed by the current convention.
Thus matter gains precisely the four-force lost by the field.
\end{solution}

\begin{solution}
Conservation gives
\[
k^\nu\nabla_\mu(\rho k^\mu)
+\rho k^\mu\nabla_\mu k^\nu=0.
\]
The second term must be proportional to \(k^\nu\), because projection
transverse to \(k^\nu\) leaves no other term.  Hence
\(k^\mu\nabla_\mu k^\nu=\kappa k^\nu\).  Rescaling \(k^\mu\) sets
\(\kappa=0\); then the remaining equation is
\(\nabla_\mu(\rho k^\mu)=0\), conservation of ray number.
\end{solution}

\begin{solution}
Write \(n=n_0+\delta n\), \(\rho=\rho_0+\delta\rho\), and
\(u^\mu=(1,\delta\bm v)\).  To first order,
\[
\partial_t\delta n+n_0\bm\nabla\cdot\delta\bm v=0,\qquad
(\rho_0+p_0)\partial_t\delta\bm v+\bm\nabla\delta p=0.
\]
Using the adiabatic relation
\(\delta p=c_s^2\delta\rho\) and the first law yields
\[
(\partial_t^2-c_s^2\nabla^2)\delta\rho=0,\qquad
c_s^2=\frac{\dd p}{\dd\rho}.
\]
Causality requires \(0\leq c_s^2\leq1\).
\end{solution}

\begin{solution}
The isentropic Euler equation may be written
\[
u^\mu\nabla_\mu(hu_\nu)+\partial_\nu h=0.
\]
Contract with \(\xi^\nu\).  Stationarity gives
\(\Lie_\xi h=0\).  The Killing equation also gives
\[
hu^\mu u^\nu\nabla_\mu\xi_\nu=0.
\]
Therefore
\[
u^\mu\nabla_\mu(h\,u_\nu\xi^\nu)=0,
\]
so the quantity is constant on a streamline.
\end{solution}

\begin{solution}
The collisionless distribution obeys
\[
p^\mu\partial_\mu f
-\Gamma^i{}_{\alpha\beta}p^\alpha p^\beta
\frac{\partial f}{\partial p^i}=0,
\]
which is \(\dd f/\dd\lambda=0\) along a null geodesic.  Since
\(f\propto I_\nu/\nu^3\), that ratio is conserved.  Integrating over
frequency gives \(I_{\rm obs}=I_{\rm em}(\nu_{\rm obs}/\nu_{\rm em})^4\);
for cosmological redshift this is the \((1+z)^{-4}\) dimming law.
\end{solution}

\begin{solution}
Take
\(A_\mu=\Re[(a_\mu+\cdots)e^{\ii\theta/\epsilon}]\) and
\(k_\mu=\partial_\mu\theta\).  The leading Maxwell equations give
\(k^2=0\) and \(k\cdot a=0\).  At the next order,
\[
k^\nu\nabla_\nu a^\mu
=-\frac12(\nabla_\nu k^\nu)a^\mu+\alpha k^\mu.
\]
Separating amplitude \(a^\mu=\mathcal A e^\mu\) leaves
\(k^\nu\nabla_\nu e^\mu\propto k^\mu\); the proportional part is gauge, so
physical polarization is parallel transported.
\end{solution}

\begin{solution}
\[
T_{\mu\nu}
=\nabla_\mu\phi\nabla_\nu\phi
-\frac12g_{\mu\nu}(\nabla\phi)^2-g_{\mu\nu}V.
\]
For null \(\ell^\mu\),
\(T_{\mu\nu}\ell^\mu\ell^\nu=(\ell\cdot\nabla\phi)^2\geq0\), so the null
condition always holds.  The weak condition holds for nonnegative \(V\).
The strong-condition contraction contains
\((v\cdot\nabla\phi)^2-V\) in four dimensions, so sufficiently large positive
\(V\) violates it and permits accelerated expansion.
\end{solution}

\begin{solution}
Metric variation gives
\[
T^{\mu\nu}(x)
=\frac{m}{\sqrt{-g(x)}}\int
u^\mu u^\nu\,\delta^{(4)}(x-z(\tau))\,\dd\tau.
\]
Testing \(\nabla_\mu T^{\mu\nu}\) against a smooth covector and integrating by
parts along the worldline produces
\(-m\int (Du^\nu/D\tau)\,\psi_\nu\,\dd\tau\).  Conservation for arbitrary
\(\psi_\nu\) is therefore equivalent to \(Du^\nu/D\tau=0\).
\end{solution}

\begin{solution}
In \(d\) dimensions,
\[
T^\mu{}_\mu
=F^{\mu\rho}F_{\mu\rho}
-\frac d4F_{\rho\sigma}F^{\rho\sigma}
=\left(1-\frac d4\right)F^2.
\]
Under \(g_{\mu\nu}\to e^{2\omega}g_{\mu\nu}\), the measure contributes
\(e^{d\omega}\) and the two inverse metrics in \(F^2\) contribute
\(e^{-4\omega}\).  The action and trace are conformally invariant only for
\(d=4\).
\end{solution}
